% !TEX root = ../main.tex

\chapter{基于 Jetson Orin 与 Lite3 的视觉目标导航真机部署与实验验证}

\section{问题分析}

第三章已经在 MuJoCo 中验证了分层闭环系统的可运行性，但仿真闭环通过并不意味着真机部署自然成立。仿真中的图像由渲染器稳定生成，平台接口由仿真程序直接控制，地面接触、通信时延和安全边界也相对可控；真实 Lite3 则需要在有限算力、真实相机、无线网络、运动主机和实验场地约束下持续运行。因此，本章面对的问题是：如何将已经在仿真中验证过的动作扩散视觉导航系统迁移到 Jetson Orin 与 Lite3 真机上，并验证其能否完成真实场景下的视觉目标导航与探索任务。

真机部署首先受到边缘计算条件限制。高层 NoMaD 推理必须在 Orin 上完成模型加载、图像预处理、扩散采样和航点输出；若扩散采样耗时过长，闭环频率将下降，机器人会在两次高层决策之间执行过旧的速度参考。其次，真实视觉输入与仿真图像存在分布差异，光照变化、运动模糊、动态遮挡和相机曝光都会影响距离预测和轨迹生成。再次，高层航点不能直接作用于 Lite3，必须经过 PD 映射、速度限幅、UDP 桥接、心跳维护和后台速度刷新，才能转化为运动主机可连续接收的底层速度通道。上述任一环节不稳定，都可能使离线算法优势无法转化为真实运动能力。

除算法和算力之外，真机实验还包含仿真中不突出的工程问题。真实拓扑地图需要现场采集和组织，导航过程中需要实时选择子目标；系统需要同时支持导航、探索、可视化、无可视化、图像保存和日志记录；实验场地、电池续航、关节温度、人工监护和急停手段也会限制实验组织方式。由此，本章将重点放在真机双链路架构、Orin 分阶段联调、真实拓扑地图与滑动窗口导航、安全机制以及多场景实验结果上，分析第二章统一推理模块的算法配置和第三章系统架构在真实平台上的实际表现。

\section{真机系统双链路与 Orin 部署架构}

\begin{figure}[htbp]
    \centering
    \includegraphics[width=0.8\linewidth]{shuanglianlu.png}
    \caption{Orin 与 Lite3 真机部署总体架构图}
    \label{fig:shuanglianlu}
\end{figure}

为了把仿真阶段形成的视觉导航闭环迁移到 Lite3 真机，本文并不是简单地把 NoMaD 模型放到 Orin 上运行，而是将系统实现整理为两条相互交叉的链路。第一条是\textbf{横向实时推理链路}，描述数据如何在一个控制周期内从相机图像流动到速度指令；第二条是\textbf{纵向分层控制链路}，描述不同软件层和硬件层如何分工协作。前者实现一帧图像变成一条 twist 指令，后者实现这条指令安全、连续地进入 Lite3 运控。

\subsection{横向实时推理链路：从图像到 twist 指令}

横向链路是每个高层导航周期内实际发生的数据流，可概括为
\begin{equation}
\begin{aligned}
\text{Camera Images} &\rightarrow \text{Visual Encoder}
\rightarrow \text{Distance Predictor} \\
&\rightarrow \text{Diffusion Policy}
\rightarrow \text{PD Mapping}
\rightarrow \text{Twist Command}.
\end{aligned}
\end{equation}
式中，Camera Images 表示相机图像输入，Visual Encoder 表示视觉编码器，Distance Predictor 表示距离预测器，Diffusion Policy 表示扩散策略头，PD Mapping 表示比例-微分控制映射，Twist Command 表示包含线速度和角速度的速度指令；箭头表示一个高层控制周期内的数据流向。其中，相机图像首先经过统一预处理，形成 NoMaD 所需的历史观察张量。本文部署中，模型每次接收当前帧及其之前三帧，共四帧第一视角图像；导航模式还会输入一张实时子目标图像，并关闭 goal mask，使目标 token 真实参与条件编码。

视觉编码器是横向链路的第一段核心模型。它把四帧观察图像编码为观察 token，并把当前帧与目标图像拼接后编码为目标 token，再通过 Transformer 融合为统一条件特征。这个条件特征并不是只供扩散模型使用，而是同时喂给距离预测器和扩散策略头，因此视觉编码器的输出质量会同时影响局部定位和局部轨迹生成。

距离预测器位于视觉编码器之后，承担拓扑地图局部定位功能。真机导航时，系统不会直接把最终目标图像作为唯一目标，而是在当前定位指针附近构造滑动窗口，将窗口内每个候选节点分别与当前观察组成条件输入，预测当前状态到候选节点的离散距离。距离最小的候选节点用于更新定位指针，并进一步确定本轮扩散策略实际追踪的实时子目标。

扩散模型位于距离预测器之后，负责在实时子目标条件下生成未来局部航点序列。本文真机部署中采用 DDIM/DDPM 形式的轨迹去噪采样，输出的是机器人局部坐标系下的多步航点，而不是底层关节动作。随后，中层 PD 映射根据选定航点计算前向速度、横向速度和偏航角速度，形成 Lite3 运动主机可接受的 twist 指令。横向链路由此完成从视觉输入到速度参考的端到端闭环。

\subsection{纵向分层控制链路：从 NoMaD 到运控}

纵向链路体现真机系统的工程分层，可概括为
\begin{equation}
\begin{aligned}
\text{High-level NoMaD Host}
&\rightarrow \text{Middle-level State Machine and PD Control} \\
&\rightarrow \text{Low-level UDP Bridge}
\rightarrow \text{Lite3 Locomotion Controller}.
\end{aligned}
\end{equation}
式中，High-level NoMaD Host 表示高层 NoMaD 导航主机，Middle-level State Machine and PD Control 表示中层状态机与 PD 控制映射，Low-level UDP Bridge 表示底层 UDP 通信桥接，Lite3 Locomotion Controller 表示 Lite3 运动控制主机；箭头表示控制命令从高层决策到真实运动执行的传递路径。最高层是运行在 Jetson Orin 上的 NoMaD 导航主机，负责相机读取、拓扑地图管理、模型推理、任务队列、可视化、录像和日志。该层主要处理视觉导航意义上的目标、节点、航点和任务状态，并不直接接触 Lite3 的关节级控制。

中层由状态机和 PD 控制映射组成。状态机负责组织空闲、起立、导航、探索、遥操作和急停等运行状态，决定何时进入导航、何时停止、何时触发急停或恢复；PD 映射则把扩散模型输出的局部航点转换为速度参考，并施加线速度、横向速度和角速度限幅，避免生成式策略偶发输出过大动作。

下层是真机桥接和 Lite3 运动主机。桥接层通过 UDP 与 Lite3 通讯，维护心跳和后台速度刷新线程，并将中层速度参考拆分为运动主机能够识别的底层速度通道。Lite3 自身的运控系统负责步态、平衡和关节级执行。这样，高层 NoMaD 可以保持与仿真阶段一致的输入输出接口，而真实机器人所需的通讯、起立、模式切换、心跳和安全停机都被封装在中下层。

两条链路在运行时并不是彼此独立的。横向链路在每个高层周期产生新的 twist 参考，纵向链路负责保证该参考在真实平台上被连续、安全地下发；当高层推理频率低于底层速度刷新频率时，桥接层会在相邻两次 NoMaD 推理之间持续发送最近一次限幅后的速度命令。正是这种横向端到端感知决策与纵向分层执行保护的结合，使得本文系统既能保留 NoMaD 的视觉导航能力，又能满足 Lite3 真机部署的实时性和安全性要求。

\subsection{Orin 真机部署架构}

\paragraph*{边缘计算平台的角色}

在本文部署方案中，Jetson Orin 承担高层导航主机的角色，即在其上完成相机读取、NoMaD
模型加载、扩散采样推理、状态机调度、导航主机交互和图像记录等工作。之所以选择 Orin 这类边缘计算平台，是因为它能够在机载条件下提供相对完整的 GPU
推理能力，从而使高层视觉导航不必完全依赖远程服务器。

从硬件规格来看，Jetson Orin 系列提供了多种配置等级，不同模块在 GPU 规模、统一内存容量和可配置功耗上并不完全相同。对于未在实验记录中逐项确认的核心数、Tensor Core 数量或持续功耗，本文不将其作为真机能力结论，而是把板端是否能够稳定调用 GPU 推理能力、是否能稳定加载模型、以及端到端推理频率作为部署评价依据。相比之下，桌面级 RTX 4090 更适合离线训练和统计实验，但其供电、散热和体积条件并不适合直接随 Lite3 移动部署。

与其他可选的边缘计算平台相比，Jetson Orin 的优势主要体现在 CUDA 生态兼容性上。例如，Raspberry Pi 4/5 虽然功耗更低（约 5--
15\,W），但不具备 CUDA 推理能力，无法直接运行 PyTorch GPU 模型；英特尔 NUC 系列提供了较强的 CPU 算力，但缺乏高效的 GPU
推理后端，纯 CPU 推理下 NoMaD 的运行频率将远低于实时闭环要求；Google Coral 等 TPU 加速器虽然能效比高，但对 PyTorch
模型的支持需要额外的模型转换步骤，且对扩散模型的迭代采样支持不够成熟。综合考虑推理能力、功耗、生态兼容性和机载可行性，Jetson Orin
是当前足式机器人视觉导航部署中较为合理的选择。

但与此同时，Orin 也带来了明确约束。Orin 的算力与显存有限，因此高层推理、图像显示、图像保
存和控制桥接必须争夺同一套本地资源。这意味着真机部署必须从一开始就区分``调试模式''和``正式运行模式''：调试时允许开启可视
化和更多日志，正式运行时则优先保障闭环控制稳定性。

\paragraph*{真机执行链路与资源约束}

在双链路设计的基础上，Orin 真机部署进一步落实为可执行的软件链路：相机图像进入高层 NoMaD 推理模块，模型输出航点后交给中层航点到速度命令映射，随后通过 Lite3 UDP 控制桥接进入运动主机。这一工程链路可以概括为

\begin{equation}
\text{Camera} \rightarrow \text{NoMaD Inference} \rightarrow \text{航点到速度命令映射} \rightarrow \text{UDP Bridge} \rightarrow \text{Lite3}.
\end{equation}
式中，Camera 表示相机输入，NoMaD Inference 表示高层模型推理，“航点到速度命令映射”表示将航点转换为速度命令的中层映射，UDP Bridge 表示真机通信桥接，Lite3 表示最终接收命令的机器人平台。

需要进一步说明的是，真机相机并不要求原始输出分辨率与视觉编码器输入尺寸完全一致。本文部署链路会在推理前自动将图像缩放到策略所需的输入尺寸，例如不同视觉骨干可能对应 $96\times96$ 或 $98\times98$ 输入。因此，对真机部署更关键的要求不是相机必须直接输出编码器尺寸，而是预处理过程与训练配置保持一致。在此基础上，相机标定可作为增强步骤，而非本文实验开展的必要前置条件：本文所使用的相机并未经过标定。


图像预处理流水线是连接相机原始输出与视觉编码器输入的关键环节。本文采用的预处理步骤包括：（1）从 USB 相机读取 $640\times480$ 的 BGR
图像；（2）将 BGR 格式转换为 RGB 格式；（3）将图像缩放至策略配置指定的输入尺寸（如 $96\times96$ 或
$98\times98$）；（4）将像素值归一化到 $[0,1]$ 范围并按照 ImageNet
均值和标准差进行标准化处理，使其与训练阶段的数据增强流水线保持一致。整个预处理过程在 CPU 上完成，耗时通常不超过 2\,ms，相对于推理耗时可忽略不计。

在推理效率优化方面，Orin 平台还具备引入 NVIDIA TensorRT 高性能推理引擎的潜力。TensorRT 是 NVIDIA
提供的面向深度学习模型部署的高性能推理引擎，能够通过算子融合、精度降低（FP16/INT8）和内存优化等手段显著提升推理速度。对于 NoMaD 的视觉编码器部分，理论上可以通过模型转换流程进一步压缩视觉编码耗时。然而，扩散采样模块由于其迭代特性和动态控制流，TensorRT 转换的难度较高，目前本文暂未实施这一优化。未来若需要在更高控制频率下运行，TensorRT
优化将是一个值得探索的方向。

在真机部署阶段，高层 NoMaD 仍保持与仿真阶段一致的策略输入输出定义，而低层机器人通讯和步态细节则被吸收到桥接层与运动主机之中。从系统工程角度看，并不是另起一套真机专用算法，而是在仿真闭环系统的基础上切换到真实后端。

\paragraph*{真机桥接层}

Lite3 的运动主机并不直接接受扩散轨迹或视觉目标图像这类高层抽象概念，它接受的是速度指令、心跳信息和模式切换命令。因此，在 Orin
和机器人之间必须存在一个真机桥接层，用来完成三项关键任务：其一，把高层 motion command 映射为底层速度控制命令；其二，维护与运动主机之间的 UDP
通讯与心跳；其三，对外暴露统一平台接口，使导航主机与状态机能够在不理解底层协议细节的前提下驱动真机。

在 UDP 通讯链路中，桥接层并不向运动主机传递扩散轨迹或视觉目标，而是将中层生成的速度参考转化为 Lite3 可接收的速度通道，并维持必要的心跳与速度刷新。本文设置心跳频率为 $5\,\text{Hz}$，速度刷新频率为 $25\,\text{Hz}$；由于高层 NoMaD 推理频率低于速度刷新频率，桥接层会在两次高层推理之间持续发送最近一次限幅后的速度参考，以满足运动主机对连续控制流的要求。

从端到端延迟预算的角度分析，完整的控制链路延迟可以分解为以下几个环节：相机图像采集约 25--40\,ms（取决于曝光时间和 USB 传输），图像预处理约
2\,ms，视觉编码约 57\,ms，扩散采样约 87\,ms，航点到速度命令映射耗时不超过 1\,ms，UDP 发送与网络传输约 1--
3\,ms。因此，从相机捕获图像到机器人收到对应速度命令的总延迟约为 170--190\,ms，对应约 5--6\,Hz 的闭环控制频率。对于当前保守低速行走场景（约 0.12\,m/s），这一延迟水平是可接受的；但若未来需要支持更高速度的运动（如 0.8\,m/s 以上），则需要通过减少 DDIM 步数或引入 TensorRT
优化来进一步压缩推理延迟。

通讯故障处理是桥接层设计中不可忽视的环节。本文桥接层采用的保护策略更偏向``上层主动检测与软急停''，而不是在桥接层内部维护一个固定发送失败计时器。具体来说，UDP 发送异常会被记录并限频告警；交互式导航主机会周期性检查摔倒状态和运动主机连接状态，当机器人摔倒或运动主机状态长时间未更新时，自动触发急停并退出当前任务。任务正常结束时，系统发送零速度并释放手动控制；任务失败或急停时，则进入软急停流程。

这层桥接设计与第三章中的平台接口是一致的。它既是``同一高层架构、双后端复用''在工程上的落地形式，也是真机部署安全性的关键保障。没有桥接层，导航主机将不得不直接操作底层通讯协议，这会使系统既难以维护，也难以在不同硬件平台之间扩展。

\paragraph*{真机部署中的时序等待机制}

真机部署中的时序等待并不是算法层面的推理暂停，而是围绕相机感知、运动主机、UDP 通讯和人机交互所建立的工程时序保护。与 MuJoCo 仿真不同，真实平台中的相机打开、机器人起立、模式切换和速度指令下发均需消耗物理时间；若上层程序在这些环节不等待或等待不足，则可能出现相机帧尚未稳定、机器人尚未站稳、运动主机尚未完成模式切换而高层策略已经开始下发速度的情况。因此，本文将这些等待机制视为真机系统可靠性的一部分，而非临时调试措施。

为便于讨论，本节先在系统层面抽象出四类接口：相机感知接口、平台运动接口、通讯桥接接口与人机交互接口。相机感知接口对外提供稳定的图像帧序列；平台运动接口对外提供起立、模式切换与停机等动作请求；通讯桥接接口将上层速度参考转化为底层连续指令流；人机交互接口处理操作者输入并维持空闲状态下的命令轮询。每一类接口都对应一组特定的时序约束，下面分别说明。

第一类等待存在于相机感知接口中。当 USB 或 CSI 相机首次打开失败时，系统以约半秒的间隔进行最多三次重试；首次成功打开后，还会丢弃数帧作为预热。该设计的目的是避开 USB 设备刚释放、刚插入或曝光自动调整尚未稳定时的异常帧，确保进入推理模块的图像分布与训练数据保持一致。

第二类等待存在于平台运动接口的起立与模式切换流程中。当向运动主机发送起立指令后，机器人需要约三秒完成站立姿态稳定；随后向自主控制模式与移动控制模式的切换分别需要约半秒的固定等待。上层接口允许操作者通过单一参数指定起立的总等待时间，若该值大于底层默认等待，则在底层等待之外补足相应时长。本文实验中将起立总等待设置为三秒，即不在底层之外额外延长；当现场观察到机器人起立后仍有明显晃动时，可将该参数适当增大以延长稳定时间。

第三类等待用于维持实时控制节拍。Lite3 通讯协议要求速度轴指令持续下发，因此通讯桥接接口在后台维持两条独立线程：速度刷新线程以约 25\,Hz 的频率重复下发最近一次速度参考，对应约 40\,ms 的循环间隔；心跳线程以 5\,Hz 的频率发送心跳信号，对应约 200\,ms 的循环间隔。这两个频率均满足运动主机对底层速度命令和心跳频率的最低要求。由于高层 NoMaD 推理频率约为 6--7\,Hz，明显低于底层速度刷新频率，桥接接口会在相邻两次高层航点更新之间持续重复最近一次限幅后的速度命令，使运动主机看到连续控制流而非稀疏突发指令。

第四类等待存在于人机交互接口的空闲与遥操作状态中。当系统处于空闲状态等待新命令时，导航主机以约 50\,ms 的周期轮询终端输入并保持站立姿态；键盘遥操作模式采用同样的循环间隔，对应约 20\,Hz 的手动控制节拍。该等待量在保证主机能够及时响应起立、导航、探索与急停等在线命令的同时，避免空闲状态下 CPU 资源被无意义循环占满。对于真机安全而言，空闲循环中的短等待还与连接陈旧检测协同工作：当运动主机状态超过两秒未更新时，导航主机会触发自动退出或急停处理。

\section{真机部署流程与联调结果}

\subsection{阶段一：Orin 端独立调试}

真机部署的第一阶段为 Orin 端独立调试，即在 Lite3 平台未接入的条件下，先验证相机读取、模型加载、CUDA 环境、端到端推理频率以及图像保存链路是否满足部署要求。该阶段的目的在于将高层推理子系统与机器人控制子系统的故障域加以隔离：当高层推理本身存在异常时，过早接入真机会使两个子系统的问题相互掩盖，进而显著增加故障定位的复杂度。

阶段一中需要依次验证的内容包括：相机能否稳定提供图像帧；NoMaD 模型权重能否在 Orin 上完成加载并被正确分配到 GPU；统一推理模块能否在板端输出合理的航点；DDIM 等参数下的端到端推理频率是否处于可用区间；以及实验记录能否完整保存。从研究意义上看，该阶段不仅验证了软件链路的连通性，更确认了第二章形成的统一推理模块在边缘计算平台上的可迁移性。

相机标定的处理方式同样在阶段一中确定。当相机原始画面中的直线结构基本正常、边缘畸变较小时，可直接进入独立调试，因为部署链路会自动完成输入尺寸对齐；当相机存在明显桶形畸变或针孔模型偏差时，则需先通过离线标定获得相机内参与畸变参数，再在 Orin 端启用桥接层去畸变。考虑到部署可用性与工程可维护性的平衡，本文将相机标定实现为可选启用的桥接层功能，而非所有实验场景的强制前置步骤。

除畸变问题外，真机相机的宽高比同样需要单独说明。本文 USB 相机的原始输出分辨率为 $640\times480$，而 NoMaD 基线视觉编码器的输入尺寸为 $96\times96$。直接将 $4:3$ 图像缩放为 $1:1$ 会引入一定横向几何压缩；但从学习系统的角度看，只要训练阶段、拓扑地图节点、目标图像和实时观测均采用一致的预处理方式，则该变换在输入分布一致性上仍然成立。本文首轮真机部署采用统一缩放策略以保留完整视场，并在后续实验中保留对比例保持或裁剪策略进行对比的可能。

\subsection{阶段一实测结果与分析}


\begin{table}[htbp]
    \centering
    \caption{Orin 阶段一独立调试实测结果}
    \label{tab:orin-stage1-results}
    \begin{tabular}{p{0.20\linewidth}p{0.24\linewidth}p{0.42\linewidth}}
        \toprule
        测试项 & 实测结果 & 结果说明 \\
        \midrule
        相机读取 & $640\times480$，约 $40.0\,\mathrm{ms}$，约 25.0 FPS & 说明 USB
        相机链路稳定，能够持续提供导航输入图像。 \\
        模型加载 & 2.87\,s，GPU 推理设备正常识别 & 说明 Jetson Orin 集成 GPU 已被正确识别，NoMaD
        权重可在板端完成加载。 \\
        基准推理（DDIM-5） & 视觉编码 56.9\,ms；扩散采样 86.7\,ms；端到端 143.6\,ms；约 7.0\,Hz &
        说明高层推理模块在 Orin 上已达到可用于后续桥接验证的运行频率。 \\
        完整流水线 & 首轮 2895.5\,ms；后续约 151--157\,ms & 说明系统存在明显冷启动开销，但稳态运行阶段可保持约 6.4--
        6.6\,Hz。 \\
        桥接层相机接口测试 & 连续 30 帧均成功读取 & 说明桥接层相机包装与独立相机测试结果一致。 \\
        部署资产完整性检查 & 模型、控制器与实验资产齐备 & 说明阶段二联调所需的模型、主机、控制器和资产链路完整。 \\
        \bottomrule
    \end{tabular}
\end{table}

从结果上看，阶段一最重要的结论有三点。第一，相机链路与模型链路已经分别通过验证。第二，推理基准测试给出的约 7.0\,Hz
更适合作为 Orin 高层推理能力的代表值，因为它排除了相机初始化和首轮图优化带来的额外开销；相比之下，端到端流水线测试的原始平均值
427.3\,ms 虽然真实，但被第一轮 2895.5\,ms 的冷启动延迟显著拉高，不能直接代表系统进入稳态后的闭环能力。第三，完整流水线在首轮之后基本稳定在 151
--157\,ms 区间，说明当前 DDIM-5 配置下系统已经具备进入 Lite3 低速联调的实时性基础。

相机读取的 FPS 为 25，意味着图像采集不会成为推理链路的瓶颈，因为高层推理频率仅约
7\,Hz，相机能够以远高于推理需求的速率供给图像帧。模型加载耗时 2.87\,s 属于一次性开销，在系统启动阶段完成，不影响运行时性能；但这一数值也表明，若系统需要在运
行中动态切换模型权重（例如从导航模型切换到探索模型），则需要额外考虑模型切换的延迟成本。基准推理中视觉编码占 56.9\,ms、扩散采样占 86.7\,ms
的分解数据表明，当前推理延迟的主要贡献来自扩散采样而非视觉编码，这与第二章中的消融实验结论一致，也进一步支持了``减少 DDIM 步数是降低推理延迟最有效手段''的判断。

因此，阶段一结果表明，Orin 端高层推理子系统已经满足进入联合调试的前置条件：一是图像输入稳定，二是模型能够使用 GPU 推理能力，三是高层输出航点数值合理且非零，四是部署资产链路完整。

\subsection{阶段二：Orin 与 Lite3 联合调试}

阶段二是 Orin 与 Lite3 联合调试阶段，包括网络连通验证、UDP 桥接测试、起立测试、低速速度测试、交互式导航测试和目标导航测试。该阶段采用从低风险动作向完整闭环渐进的组织原则：足式机器人各底层环节的不稳定性会在高层导航阶段被放大，将测试粒度按风险等级排序有助于在最早的子环节中暴露故障。

在具体顺序上，首先确认 Orin 与 Lite3 运动主机之间的网络与心跳是否正常；其次执行起立测试，验证机器人能否从趴卧状态平稳过渡到自主控制模式；随后进行低速 twist 测试，验证桥接层能否在固定低速下连续下发速度命令；最后再启动导航主机驱动的探索或目标导航流程。

阶段二各子步骤的具体内容如下。网络连通测试要求 Orin 与 Lite3 运动主机处于同一局域网内，并能够完成基本通信；心跳测试要求桥接层启动后，运动主机能够确认收到有效心跳包；起立测试通过桥接层发送起立指令，观察机器人是否能从趴卧状态平稳过渡到站立姿态，并验证控制权限是否已正确移交给 Orin 端；低速 twist 测试以小幅线速度和零角速度驱动机器人短距离前进，用于验证速度命令的方向性和幅度合理性；交互式导航测试允许操作者指定临时目标图像，观察机器人是否能朝目标方向产生合理运动响应。

阶段二结果表明，Orin 与 Lite3 的联合控制链路已具备进一步闭环验证的基础。

\section{真实拓扑地图采集与滑动窗口导航}

本节把本文构建的真机导航系统作为一个完整流程进行整理。与单张目标图像导航不同，本文系统在真实场景中采用\textbf{先采集有序拓扑地图，再沿拓扑地图滑动推进}的方式完成目标导航。其核心思想是：先让操作者沿真实路线拍摄一串第一视角图像，形成从起点到终点的视觉路标序列；正式运行时，Lite3 从序列起点附近出发，系统在局部滑动窗口中用距离预测器估计当前最接近的节点，再把由该节点推进得到的实时子目标图像送入扩散策略，生成下一段局部航点。机器人不断重复这一过程，直到定位指针推进到指定目标节点。

\subsection{导航输入与实时子目标}

真机端的高层算法仍然由 NoMaD 模型承担，但真实导航任务的输入组织与单张目标图像导航不同。系统并不是把终点图像一次性作为固定目标输入，而是根据真实拓扑地图和滑动窗口机制，在每个高层周期选择一个实时子目标图像。

相机与预处理模块负责从 Orin 连接的 USB/CSI 相机中读取实时图像，并将图像转换为与训练配置一致的输入张量。每次高层推理使用当前帧及其之前三帧，共四帧历史观察图像；若系统启动初期历史帧缓存不足，系统会重复最早帧补齐上下文。导航模式下还会额外输入一张目标图像，但这张图并不是整段任务的最终终点图，而是本轮滑动窗口机制选出的实时子目标节点图像。此时视觉掩码关闭，目标 token 参与视觉条件编码；探索模式则使用被掩盖的目标条件。

视觉编码器负责把上述四帧观察图像和实时子目标图像编码为统一条件特征。在 NoMaD 的 ViNT 结构中，四帧观察图像分别经过观察编码器得到观察 token，当前帧与目标图像沿通道维度拼接后进入目标编码器得到目标 token，随后这些 token 经过 Transformer 自注意力融合，并池化为固定维度的条件向量。这个条件向量被后续两个头共享，因此视觉编码器、距离预测器和扩散策略头不是彼此独立的三段程序，而是共享同一个视觉条件空间。

距离预测器用于拓扑地图局部定位。给定当前观察和滑动窗口中的若干候选节点图像，系统会把同一组观察帧分别与窗口内每一张候选图像组成输入，预测当前状态到候选节点的离散距离。距离最小的候选节点被视为当前视觉上最接近的位置，再经过反跳变限制更新定位指针。这一距离不是物理米制距离，而是训练数据中相对导航步数意义下的距离估计。

扩散策略头负责生成局部航点。系统不会把整个拓扑地图或最终终点一次性送入扩散模型，而是只把滑动窗口选出的一个实时子目标节点作为目标条件。扩散策略在该条件特征上通过 DDIM/DDPM 去噪生成未来若干步局部航点，随后中层控制器从预测轨迹中选取指定索引的航点，将其转换为 Lite3 可执行的前向速度、横向速度和偏航角速度。最后，真机桥接层通过 UDP 将速度命令持续发送给 Lite3 运动主机，并在两次高层推理之间保持最近一次速度参考。

\subsection{真实拓扑地图的采集与组织}

在正式导航之前，必须先在真实环境中采集拓扑地图。采集阶段通过 Orin 端相机接口沿路线记录图像节点。

采集方式分为自动和手动两种。自动模式按照设置的时间间隔保存图像，适合操作者能够以稳定速度沿路线移动相机的场景；手动模式程序会显示实时预览窗口，操作者每到一个希望记录的节点位置后保存当前节点。

采集得到的图像按照文件名排序后构成有序序列
\begin{equation}
\mathcal{M}=\{I_0,I_1,\ldots,I_N\},
\end{equation}
其中 $\mathcal{M}$ 表示真实拓扑地图图像序列，$I_0$ 对应导航起点附近的第一视角，$I_N$ 通常对应默认终点，$N$ 表示序列最后一个节点索引。节点编号保证与真实路线的前进顺序一致，这是系统能够滑动推进的前提。

在节点间距上，过密的拓扑地图会增加存储和匹配开销，但相邻节点高度相似时不一定显著提升效果；过疏的拓扑地图则会造成相邻目标之间的视觉跳变，使模型难以在中间位置生成稳定轨迹。室内走廊拓扑地图中，相邻节点间隔距离通常为 1 m 至 1.5 m；在较为空旷的地面场景中，相邻节点间隔距离可控制在 0.5 m 至 1 m。

\subsection{滑动窗口导航机制}

真实导航开始时，任务队列将定位指针初始化为
\begin{equation}
c_0=0,
\end{equation}
式中 $c_0$ 表示初始定位指针，取值 $0$ 表示默认机器人从拓扑地图的第一张图像附近出发。目标节点 $g$ 默认为序列最后一个节点 $N$，也可以通过指定拓扑地图中的某个图像节点文件名来提前结束任务。需要强调的是，$g$ 只是当前任务的终点上界，并不表示每个推理周期都直接把 $I_g$ 作为扩散目标。

在第 $t$ 个高层周期，系统以当前定位指针 $c_t$ 为中心构造局部候选窗口。本文采用的窗口边界为
\begin{equation}
s_t=\max(c_t-r,0),\qquad e_t=\min(c_t+r+1,g),
\end{equation}
式中，$s_t$ 和 $e_t$ 分别表示第 $t$ 个高层周期中局部窗口的左、右边界索引，$c_t$ 为当前定位指针，$r$ 为窗口半径，$g$ 为目标节点索引上界。
候选集合实际为
\begin{equation}
\mathcal{W}_t=\{I_{s_t},I_{s_t+1},\ldots,I_{e_t}\}.
\end{equation}
式中，$\mathcal{W}_t$ 表示当前局部候选窗口，$I_{s_t}$ 到 $I_{e_t}$ 表示窗口内的候选拓扑地图节点图像。实际右边界使用 $c_t+r+1$，是为了在当前节点前后搜索的同时保留一个更靠前的候选节点，便于机器人接近当前节点后继续向下一站推进。起点附近窗口左端被裁剪到 0，终点附近窗口右端被裁剪到 $g$，因此窗口会先从起点附近展开，中途随 $c_t$ 向后滑动，最后停在目标节点 $g$ 之前或包含 $g$ 的局部范围内。

随后，系统把同一组观察帧分别与窗口内每个候选节点图像组成目标条件，调用视觉编码器和距离预测器得到一组距离估计：
\begin{equation}
d_i = D_\theta(O_t, I_i),\qquad i\in[s_t,e_t].
\end{equation}
式中，$d_i$ 表示当前观察到候选节点 $I_i$ 的离散距离估计，$D_\theta$ 表示参数为 $\theta$ 的距离预测器，$O_t$ 表示第 $t$ 个周期的历史观察，$i$ 为窗口内候选节点索引。
距离最小的节点
\begin{equation}
\hat{c}_t=\arg\min_{i\in[s_t,e_t]}d_i
\end{equation}
式中，$\hat{c}_t$ 表示局部窗口内的最近节点估计，$\arg\min$ 表示选择使 $d_i$ 最小的候选索引。为了防止视觉误匹配导致定位指针突然跳远，实际实现时还会把该结果限制在 $[c_t-1,c_t+3]$ 范围内，并且不允许超过目标节点 $g$。这一限制使指针可以缓慢后退一格以修正误差，但最多只允许每轮向前推进三个节点。

扩散模型实际使用的目标节点记为 $u_t$。它并不总是窗口中距离最小的节点：如果距离预测器认为机器人已经足够接近当前最近节点，即 $d_{\min}<\tau$，系统会把扩散条件推进到该节点之后的下一张图像；否则仍使用当前最近节点。形式上可以概括为
\begin{equation}
u_t=
\begin{cases}
\min(\hat{c}_t+1,g), & d_{\min}<\tau,\\
\hat{c}_t, & d_{\min}\ge \tau.
\end{cases}
\end{equation}
式中，$u_t$ 表示第 $t$ 个周期实际送入扩散策略的目标节点索引，$d_{\min}$ 表示窗口内最小预测距离，$\tau$ 表示推进阈值，$\min(\hat{c}_t+1,g)$ 保证子目标不会超过任务终点。因此，距离预测器的作用是定位和推进指针，扩散策略的作用是朝实时子目标 $I_{u_t}$ 生成下一段局部轨迹。机器人执行该局部轨迹后，下一轮再读取新的相机帧、更新四帧上下文、重新构造窗口并重复上述过程。整个过程不是一次性规划全局轨迹，而是“观察--局部定位--选择子目标--生成航点--执行--再观察”的闭环。

由于真机实验未引入外部定位系统或运动捕捉系统，系统无法像仿真环境那样直接计算机器人当前位置与目标点之间的全局欧氏距离。因此，本文在真实平台上采用视觉拓扑意义下的到达判据。具体而言，当实时子目标节点已经达到或超过任务目标节点，且距离预测器给出的最近节点距离低于阈值时，认为机器人已到达目标区域：
\begin{equation}
u_t \ge g,\qquad d_{\min}<\tau_{\mathrm{arrive}}.
\end{equation}
式中，$\tau_{\mathrm{arrive}}$ 为真机任务完成阈值，本文实现中取 $\tau_{\mathrm{arrive}}=3.0$，与局部节点推进阈值保持一致。需要强调的是，该数值并非 $3\,\mathrm{m}$ 的物理距离，而是 NoMaD 距离预测头输出空间中的无量纲阈值；该输出反映当前观察与候选拓扑节点在训练分布中的相对导航距离，其物理尺度由真实拓扑地图的采集间隔间接决定。因此，真机实验中的“到达”表示机器人在视觉拓扑地图上已经推进到目标节点附近，并且全过程未触发急停或人工干预。

\section{导航/探索模式与安全机制}

\subsection{导航模式与探索模式}

真机端同样支持导航模式与探索模式两类高层运行方式。导航模式下，系统明确接收目标图像或目标节点，并由 NoMaD
生成朝向该目标的局部轨迹；探索模式下，系统弱化明确目标条件，使高层策略依赖当前视觉上下文和训练中学到的先验运动趋势继续前进。由于 goal mask
已经在训练阶段统一了这两种能力\cite{nomad2023}，因此真机端不需要切换完全不同的模型，只需要切换任务模式与目标输入即可。


\subsection{可视化模式与无可视化模式}

真机部署中另一个需要明确区分的问题是可视化模式与无可视化模式。可视化模式主要用于调试、录屏与实验观察，此时系统会显示实时相机画面、当前目标图像以及必要的状态信息；无可视化
模式则面向正式运行，优先保证闭环稳定与计算资源利用，不打开图形窗口，只保留日志与必要图像保存。

在 Orin 这类边缘计算平台上，可视化界面显示和图像刷新会消耗额外资源，从而影响推理周期稳定性。将``是否显示
图像''设计为可切换项，意味着系统可以在调试阶段最大化可观测性，而在正式运行阶段最大化稳定性。这也是本文在系统层特别强调``推理模块''和``部署模块''分离的原因之一。

\subsection{安全机制}

本文在部署链路中设置了多处可验证的安全入口，包括任务正常结束时的零速度停止、任务失败或急停时的软急停、导航主机对摔倒和连接陈旧状态的检测，以及实验者保留的遥控器或硬件断电手段。与此同时，桥接层持续维护底层心跳和后台速度刷新，减少高层推理周期抖动对运动主机的影响。

从层级化安全架构的角度，本文系统可以概括为四类保护。第一类是任务级停止：状态机在任务完成时发送零速度并释放手动控制。第二类是软件急停：当状态机进入急停或任务失败路径时，真机后端执行软急停。第三类是状态监测：导航主机在交互式运行中检查摔倒状态和运动主机连接状态，触发后会急停并退出当前任务。第四类是实验规程层面的人工保护：监护者应保留遥控器、物理急停或断电手段，用于覆盖软件链路之外的极端情况。

这些安全机制并不是与 NoMaD 无关的额外功能，而是高层生成式策略能够进入真机的先决条件\cite{chi2025diffusion,hwangbo2019anymal,lee2020learning}。因为生成式模型输出本身具有分布性，如果系统没有明确的速度限幅、恢复和急停通道，即使轨迹在大多数时刻合理，也可能在局
部异常时对平台造成风险。

\section{真机闭环实验结果}

本节给出 Lite3 真机平台上已经完成的四类实验结果：真实场景拓扑地图采集、多场景目标导航、多场景探索与室外零样本导航。所有实验均在保守速度限幅下进行（前向线速度上限 $0.12\,\text{m/s}$、偏航角速度上限 $0.35\,\text{rad/s}$），并在每次运行前后保存会话日志、运行视频和关键帧图像，以保证结果的可复现性。

\subsection{真实场景拓扑地图采集}

正式导航实验前，本文按前述滑动窗口导航流程依次在多个真实场景中采集了有序拓扑地图节点序列。采集完成后，每条拓扑地图均通过人工回放检查序列的空间连续性与朝向一致性，确保进入正式导航实验时不会因节点排序异常导致定位指针推进错误。

\begin{figure}[htbp]
    \centering
    \includegraphics[width=0.8\linewidth]{map.png}
    \caption{真实场景拓扑地图采集}
    \label{fig:map}
\end{figure}

\subsection{多场景目标导航实验}

为评估第二章统一推理模块筛选出的推理配置在真机上的实际表现，本文在前述真实拓扑地图上组织了多场景目标导航实验，并对四组推理配置进行了闭环对比。四组配置分别为：基线 DDPM 配置（沿用 NoMaD 默认采样器与步数）、最快方案（DDIM-2、CFG=0、TTS-8）、最优方案（DDIM-2、CFG=2、TTS-8）和对比方案（DDIM-3、CFG=0、无 TTS）。每组配置在多场景中各运行若干次，并以前述视觉拓扑到达判据成立且未触发急停或人工干预作为成功条件。

实验结果如表~\ref{tab:real-multiscene-navigate} 所示。在当前真实场景规模与速度限幅下，每组配置进行五次实验，四组配置在多场景目标导航中均达到 100\% 成功率。在此基础上，更具区分意义的指标是闭环频率，反映各推理配置在真实部署条件下的实时性差异。

\begin{table}[htbp]
    \centering
    \caption{多场景目标导航的真机闭环实验结果}
    \label{tab:real-multiscene-navigate}
    \begin{tabular}{lcc}
        \toprule
        推理配置 & 成功率 & 闭环频率/Hz \\
        \midrule
        基线 DDPM & 100\% & 2.89 \\
        最快方案（DDIM-2, CFG=0, TTS-8） & 100\% & 6.80 \\
        最优方案（DDIM-2, CFG=2, TTS-8） & 100\% & 3.43 \\
        对比方案（DDIM-3, CFG=0, 无 TTS） & 100\% & 5.90 \\
        \bottomrule
    \end{tabular}
\end{table}

从成功率维度看，四组配置在多场景目标导航中均达到 100\%。这一结果表明，在当前真实场景规模与保守速度限幅下，系统具有充分的鲁棒性裕度：即便最慢的基线 DDPM 配置（2.89\,Hz、约 346\,ms 单帧周期）也能稳定完成沿拓扑地图推进的导航，CFG 与 TTS 等推理质量增强项的引入也未导致额外失败。因此，在当前实验范围内，成功率指标本身难以进一步区分四组配置的性能差异，更具区分意义的指标转向了闭环频率。

闭环频率与推理配置之间呈现出与第二章离线分析高度一致的趋势。基线 DDPM 使用完整 10 步去噪，单次扩散采样耗时最长，闭环频率最低（2.89\,Hz）；将采样器替换为 DDIM 后，去噪步数大幅减少：DDIM-3 的对比方案提升至 5.90\,Hz，DDIM-2 的最快方案进一步提升至 6.80\,Hz。这一``DDIM 步数减少 $\Rightarrow$ 闭环频率单调提升''的关系直接验证了第二章关于``少步 DDIM 是降低推理延迟最有效手段''的结论。CFG 的开销同样可在频率数据中清晰读出：最优方案与最快方案除 CFG 权重外其他配置完全一致，但最优方案（CFG=2）的 3.43\,Hz 约为最快方案（CFG=0）的 6.80\,Hz 的一半，原因在于 CFG 每个去噪步需执行条件与无条件两次前向传播，等效将扩散采样计算量翻倍，这与第二章和第三章中``CFG 在提升条件牵引力的同时引入约一倍的采样开销''的结论相吻合。TTS 的开销则相对可控：最快方案（DDIM-2+TTS-8，6.80\,Hz）的频率仍高于对比方案（DDIM-3 无 TTS，5.90\,Hz），表明 TTS 候选轨迹通过批量并行采样后，其相对单条轨迹的额外耗时主要落在轻量评分函数上，远低于扩散去噪本身，这与第二章关于``批量 TTS 在 GPU 上具有较好并行效率''的结论保持一致。综合来看，真机闭环频率与第二章离线消融实验给出的相对快慢顺序基本一致，说明算法层结论在迁移到真实部署条件后未出现相反趋势。

值得进一步分析的是，尽管基线 DDPM 仅以 2.89\,Hz 运行，机器人仍然能够稳定完成所有真机导航任务。其一，本文采用的保守速度限幅（前向 $0.12\,\text{m/s}$）使得机器人在两次相邻高层推理之间至多移动约 $0.12/2.89\approx 4.2\,\text{cm}$，远小于拓扑地图的节点间距，因此低频高层更新不会在子目标层面引起明显的定位指针漂移。其二，通讯桥接层以 25\,Hz 的频率持续重复最近一次速度命令，使运动主机在两次 NoMaD 推理之间仍能看到连续控制流，避免了高层频率不足导致的速度指令稀疏问题。其三，滑动窗口拓扑地图机制为高层导航提供了空间锚点，即使单次扩散采样存在轻微抖动，距离预测器在下一周期仍可将定位指针拉回正确节点。三者共同保证了高层频率的下限被推至 $\sim 3\,\text{Hz}$ 量级仍不破坏闭环可行性。

最后，最快方案在保持 100\% 成功率的同时将闭环频率从基线的 2.89\,Hz 提升至 6.80\,Hz，约为基线的 2.4 倍。这一结果具有双重意义：一方面，基线 DDPM 的成功部署表明系统对低频高层推理具有充分容错能力；另一方面，最快方案带来的频率余量为后续支持更高速度运动和更复杂动态环境提供了方案。

需要说明的是，上述``100\% 成功率''是当前室内实验场景规模下的结论。受训练集、电池续航、场地尺寸与人工监护成本的限制，首先本文算法训练只用了室内数据集，其次本文真机实验的样本数有限；若将场景扩大到室外场景、含强光变化、动态障碍或更长走廊的环境中，不同推理配置之间的差异有望更显著，高频配置在面对快速变化场景时的优势也可能进一步放大。

\begin{figure}[htbp]
    \centering
    \includegraphics[width=0.5\linewidth]{real-robot-eg.png}
    \caption{真机目标导航过程}
    \label{fig:real-multiscene-navigate-placeholder}
\end{figure}

\subsection{多场景探索实验}

探索模式实验同样在多个真实场景中开展。与目标导航不同，探索任务不提供目标 token，而是依靠训练阶段习得的运动先验在当前视觉上下文中持续生成可行轨迹。本文设置该实验的目的并不是对探索能力进行完整定量评价，而是验证无目标条件链路在真机平台上是否能够正常运行，包括 goal mask 是否能够正确进入推理模块、扩散策略是否能够在无目标条件下输出非零且连续的航点、中层控制映射是否能够将其转化为安全速度参考，以及状态机是否能够在探索、停止和急停之间正常切换。

需要说明的是，探索任务缺少固定终点和唯一参考路径，其运行效果与场地边界、人工停止时机和安全监护策略密切相关。如果直接使用成功率、到达时间或轨迹误差等指标进行统计，容易把场地选择和人工策略的影响混入模型能力评价。因此，本文没有为探索模式单独设计大规模定量实验，而是将其定位为真机可用性验证。实验过程中，机器人能够沿走廊或开阔区域稳定前进，在接近墙壁或障碍物前完成减速，并由扩散策略生成绕行或转弯的局部航点。这一行为模式与第二章关于 goal mask 的分析一致，说明探索模式在真机平台上得到了行为层面的验证。

\subsection{室外零样本导航实验}

前述真机实验主要位于室内场景，而本文用于训练和离线评估的视觉导航数据同样来自室内路线。为了进一步检验系统在分布外环境中的可用边界，本文将 Lite3 放置到室外环境中进行零样本导航挑战。这里的``零样本''指高层 NoMaD 模型没有使用室外数据进行重新训练或微调；室外拓扑地图仅作为任务输入提供视觉目标序列，不参与模型参数更新。

室外场景相对于室内场景具有更强的视觉分布差异。首先，室外光照会受到太阳直射、阴影和曝光变化影响，同一路径在不同朝向下可能呈现明显亮度差异。其次，室外地面纹理、植被、天空、建筑外立面等视觉特征与室内走廊差异较大，容易影响视觉编码器对当前位置和目标图像关系的表征。再次，开阔区域缺少走廊墙面这类稳定几何边界，动态行人、车辆或风吹动的物体也会增加图像中的非静态因素。上述变化都会干扰视觉编码器、距离预测器和扩散策略头之间的条件分布，使室外导航成为比室内导航更强的泛化挑战。

该实验的研究目的不是证明本文系统已经具备可靠的室外自主导航能力，而是进行一次分布外压力测试：在不引入室外训练数据、不改变模型权重的前提下，观察拓扑地图滑动窗口机制和动作扩散策略能否保持基本闭环，识别视觉编码器在强光照、复杂纹理和开阔空间下可能出现的失效形式，并为后续室外数据采集和编码器重训提供依据。实验仍采用保守速度限幅和人工监护，重点观察机器人是否能够沿室外拓扑地图产生连续运动响应、定位指针是否出现明显跳变，以及是否需要人工急停介入。

\begin{figure}[htbp]
    \centering
    \includegraphics[width=0.8\linewidth]{shiwai.png}
    \caption{室外零样本导航过程}
    \label{fig:outdoor-zero-shot-navigation}
\end{figure}

从实验结果看，机器人在室外简单直线斜坡场景中完成了导航任务，但在草坪这类特征复杂且光照强烈的场景中出现失败。该结果初步表明，在未经室外数据训练的条件下，室外零样本导航更适合在视觉变化较小、路线结构较简单的场景中维持基本闭环。当室外路径具有相对连续的地面纹理、光照变化较弱且拓扑地图节点间视觉差异较小时，机器人能够沿目标序列产生连续的前进和转向响应，说明本文构建的真机链路在室外环境中并非完全失效。然而，一旦场景中出现强光直射、阴影突变、地面纹理重复或大面积开阔区域，系统的导航稳定性明显下降，表现为定位指针波动、实时子目标选择不稳定以及局部航点方向性减弱等现象。

由此可以看出，本文系统在室外环境中的主要限制并不来自 Lite3 底层运动执行或 UDP 桥接链路，而是来自高层视觉表征的分布外泛化能力。由于训练与离线评估数据主要来自室内路线，视觉编码器更熟悉走廊、墙面、门框等室内结构线索；当输入转为室外光照、植被、天空和开阔地面时，历史观察与目标图像之间的相似性判断会变得不稳定，进而影响距离预测器和扩散策略头的条件输入。因此，室外实验更适合作为本文系统泛化边界的验证，而不是作为已经具备成熟室外导航能力的证明。后续若要提升室外可用性，需要补充室外数据、重新训练或适配视觉编码器，并结合更丰富的安全约束与定量评价指标开展系统实验。


\section{本章小结}

本章围绕 Orin 与 Lite3 的真机部署问题，对统一闭环系统在真实平台上的实现与验证进行了系统阐述。

首先，本章给出了真机系统的双链路结构。横向实时推理链路把每个高层周期内``相机图像 $\to$ 视觉编码 $\to$ 距离预测 $\to$ 扩散策略 $\to$ PD 映射 $\to$ twist 指令''的端到端数据流刻画为一段从感知到执行的完整闭环；纵向分层控制链路则把整套软件实现拆分为高层 NoMaD 主机、中层状态机与 PD 控制、底层 UDP 桥接和 Lite3 运动主机四层职责，使 Orin 上的高层视觉策略能够通过桥接层与 Lite3 形成稳定的指令交互，而无需直接面对底层关节级控制。

其次，本章详细整理了 Orin 端的部署架构和真机桥接层的设计，明确了边缘计算平台的角色、相机预处理流水线和速度刷新与心跳维护机制。在此基础上，本章按``阶段一独立调试 $\to$ 阶段二联合调试''的两段式流程对真机部署进行了渐进式验证：阶段一在 Orin 端独立确认了相机链路、模型加载、CUDA 推理频率和图像保存通路，证明高层推理子系统已经具备部署条件；阶段二则在与 Lite3 联合的低风险动作上逐级验证了网络连通、起立、低速 twist 与目标导航流程。

第三，本章把仿真阶段中``状态机 + 导航主机''的任务组织方式扩展到了真机场景，重点说明了真实拓扑地图采集与滑动窗口导航机制。在该机制下，距离预测器负责沿拓扑地图推进定位指针，扩散策略则在每个周期内以推进得到的实时子目标为条件生成局部航点，机器人通过``观察—局部定位—选择子目标—生成航点—执行—再观察''的闭环逐步走完整段路径。本章还讨论了导航模式与探索模式、可视化模式与无可视化模式的切换原则，以及多级安全机制与异常恢复入口。

最后，本章给出了真实场景下的多场景目标导航实验结果。在保守速度限幅下，基线 DDPM、最快方案（DDIM-2 + CFG=0 + TTS-8）、最优方案（DDIM-2 + CFG=2 + TTS-8）和对比方案（DDIM-3 + CFG=0，无 TTS）四组推理配置的成功率均为 100\%，闭环频率分别为 2.89\,Hz、6.80\,Hz、3.43\,Hz 和 5.90\,Hz。该频率排序与第二章离线消融实验的相对快慢顺序完全一致，说明算法层得到的``少步 DDIM 是降低推理延迟最有效手段、CFG 会带来约一倍的扩散开销、TTS 排序近乎零代价''等结论在迁移到真机部署条件后仍然成立。多场景探索实验验证了依靠 goal mask 训练得到的探索模式能够在真机上正常运行，但本文未将其设计为独立的定量评测。室外零样本导航实验则进一步表明，室外强光照、复杂纹理和开阔空间会对视觉编码器与距离预测造成更强干扰，因此该实验更适合作为分布外压力测试和后续室外泛化研究的起点。

综合本章工作，Orin 与 Lite3 的真机部署不仅完成了从仿真闭环到真实闭环的过渡，也为第二章统一推理模块的算法选择和第三章分层闭环系统设计提供了真机层面的验证。需要说明的是，受场地、续航与人工监护成本限制，本章的成功率统计仍以当前实验场景规模为基础，更大规模的统计实验、更复杂场景下的差异化测试以及与拓扑地图在线维护相关的扩展工作，将在第五章的研究展望中进一步说明。
